\documentclass[11pt]{article}
\usepackage[a4paper,margin=25mm]{geometry}
\usepackage[no-math]{fontspec}
\setmainfont{Latin Modern Roman}
\usepackage{amsmath,amssymb,amsthm,xfrac,xspace,hyperref,graphicx,comment}
\excludecomment{DefTralics}
\providecommand{\C}{}
\newcommand{\ie}{i.e.,\xspace}
\newcommand{\eg}{e.g.,\xspace}
\newcommand{\etc}{etc.\xspace}
\newcommand{\cf}{cf.\xspace}
\newcommand{\resp}{resp.\xspace}
\newcommand{\smaller}{\small}
\newcommand{\Subsection}{\subsection}
\newcommand{\Subsubsection}{\subsubsection}
\newcommand{\skpt}{\smallskip}
\emergencystretch=3em
\input{author-preamble.tex}
\begin{document}
\begin{center}{\Large Stable item 1923}\end{center}
\paragraph{Source and attribution.}
Julien Marché, Maxime Wolff, \emph{The parallelogram identity on groups and deformations of the trivial character in $\mathrm{SL}_2(\mathbb{C})$}, Journal de l’École polytechnique — Mathématiques 7 (2020), publisher version, DOI 10.5802/jep.117. \href{https://jep.centre-mersenne.org/articles/10.5802/jep.117/}{Publisher article}. Authors' text: \href{https://creativecommons.org/licenses/by/4.0/}{CC BY 4.0}.
\paragraph{Editorial problem boundary.}
Prove only Theorem 1.4(2), with its first-obstruction argument retained as necessary core. For a finitely generated group with complex coefficients, a third-order character-algebra deformation at the trivial character forces its first coefficient to be a quadratic form of rank at most two on first homology. This is a necessary obstruction, with no sufficiency or representation-lifting assertion. The coordinate-algebra definitions, trace/augmentation identities, and both complete obstruction arguments are retained. Surrounding statements are context only.
\paragraph{Difficulty (editorial): H2.}
A noncommutative fourth-augmentation identity forces the cubic part of a parallelogram tangent to vanish at second order. Two distinct sixth-augmentation expansions at third order yield all three-by-three Gram minors, imposing rank at most two on the induced quadratic tangent form.
\paragraph{Editorial transcription note.}
Original author language, mathematical content and ordering are unchanged. Publisher front matter is replaced by this independent layout wrapper. Original native body and macros remain editable; bibliography is recovered from the matching cached publisher HTML. Adjacent claims are author context, not additional corpus items. Printed mathematical slips are retained without assistant repair.
\clearpage
\input{author-body.tex}
\begin{thebibliography}{99}
\bibitem{Artin} M. Artin - “On the solutions of analytic equations” , Invent. Math. 5 (1968), p. 277-291
\bibitem{Artin69} M. Artin - “On Azumaya algebras and finite dimensional representations of rings” , J. Algebra 11 (1969), p. 532-563
\bibitem{Brown} K. S. Brown - Cohomology of groups , Graduate Texts in Math. , vol. 87 , Springer-Verlag, New York, 1994
\bibitem{BrumfielHilden} G. W. Brumfiel \& H. M. Hilden - SL ( 2 ) representations of finitely presented groups , Contemporary Math. , vol. 187 , American Mathematical Society, Providence, RI, 1995
\bibitem{Chenevier} G. Chenevier - “Sur la variété des caractères p -adique du groupe de Galois absolu de \(\mathbb Q\) p ” , 2009
\bibitem{ChenevierHab} G. Chenevier - “Mémoire de HDR: Représentations galoisiennes automorphes et conséquences arithmétiques des conjectures de Langlands et Arthur” , 2013, http://gaetan.chenevier.perso.math.cnrs.fr/hdr/HDR.pdf
\bibitem{CullerShalen} M. Culler \& P. B. Shalen - “Varieties of group representations and splittings of 3 -manifolds” , Ann. of Math. (2) 117 (1983) no. 1, p. 109-146
\bibitem{Epstein} D. B. A. Epstein - “Finite presentations of groups and 3 -manifolds” , Quart. J. Math. Oxford Ser. (2) 12 (1961), p. 205-212
\bibitem{farbmargalit} B. Farb \& D. Margalit - A primer on mapping class groups , Princeton Math. Series , vol. 49 , Princeton University Press, Princeton, NJ, 2012
\bibitem{FunarMarche} L. Funar \& J. Marché - “The first Johnson subgroups act ergodically on SU 2 -character varieties” , J. Differential Geom. 95 (2013) no. 3, p. 407-418
\bibitem{goldman06} W. M. Goldman - “Mapping class group dynamics on surface group representations” , in Problems on mapping class groups and related topics , Proc. Sympos. Pure Math. , vol. 74 , American Mathematical Society, Providence, RI, 2006, p. 189-214
\bibitem{GraceYoung} J. H. Grace \& A. Young - The algebra of invariants , Cambridge Library Collection , Cambridge University Press, Cambridge, 1903
\bibitem{Liu} Q. Liu - Algebraic geometry and arithmetic curves , Oxford Graduate Texts in Math. , vol. 6 , Oxford University Press, Oxford, 2002
\bibitem{MaclachlanReid} C. Maclachlan \& A. W. Reid - The arithmetic of hyperbolic 3-manifolds , Graduate Texts in Math. , vol. 219 , Springer-Verlag, New York, 2003
\bibitem{Mukai} S. Mukai - An introduction to invariants and moduli , Cambridge Studies in Advanced Math. , vol. 81 , Cambridge University Press, Cambridge, 2003
\bibitem{GIT} D. Mumford, J. Fogarty \& F. Kirwan - Geometric invariant theory , Ergeb. Math. Grenzgeb. (2) , vol. 34 , Springer-Verlag, Berlin, 1994
\bibitem{NijenhuisRichardson} A. Nijenhuis \& R. W. Richardson - “Cohomology and deformations in graded Lie algebras” , Bull. Amer. Math. Soc. 72 (1966), p. 1-29
\bibitem{Passi} I. B. S. Passi - Group rings and their augmentation ideals , Lect. Notes in Math. , vol. 715 , Springer, Berlin, 1979
\bibitem{FriisStetkaer} P. de Place Friis \& H. Stetkær - “On the quadratic functional equation on groups” , Publ. Math. Debrecen 69 (2006) no. 1-2, p. 65-93
\bibitem{Procesi} C. Procesi - “The invariant theory of n × n matrices” , Advances in Math. 19 (1976) no. 3, p. 306-381
\bibitem{Rouquier} R. Rouquier - “Caractérisation des caractères et pseudo-caractères” , J. Algebra 180 (1996) no. 2, p. 571-586
\bibitem{Saito} K. Saito - “Character variety of representations of a finitely generated group in SL 2 ” , in Topology and Teichmüller spaces (Katinkulta, 1995) , World Sci. Publ., River Edge, NJ, 1996, p. 253-264
\bibitem{Serre} J.-P. Serre - Arbres, amalgames, SL 2 , Astérisque , vol. 46 , Société Mathématique de France, Paris, 1977
\bibitem{Stallings} J. Stallings - “Homology and central series of groups” , J. Algebra 2 (1965), p. 170-181
\bibitem{Weil} A. Weil - “Remarks on the cohomology of groups” , Ann. of Math. (2) 80 (1964), p. 149-157
\end{thebibliography}
\end{document}
